\section[Cote 156-3 : l'intermédiarité se lit sur les découpages]{Cote 156-3 : l'intermédiarité se lit sur les découpages\protect\verifie{Folder156_3}}\label{sec:156-3}

La cote 156-3, chapitre~III de \emph{Vers une géométrie des formes}, quarante pages datées des 8 et 9~juin 1986, intitulée de sa main \q{Réseaux via découpages}, propose une troisième axiomatique des réseaux : des lieux, les paires de lieux en position relative modérée, et pour chaque partie finie modérée $T$ le découpage de son complémentaire modéré $C(T)$ en composantes. La page~8 définit le \emph{tronçon ordonné} : un ensemble ordonné où, pour tout $a$, $\mathcal L_{>a}$ est filtrant décroissant et $\mathcal L_{<a}$ filtrant croissant (a, a$'$), qui est filtrant croissant et décroissant (b), divisible (c) et de cardinal $> 1$ (d) ; les 0-sphères y sont les paires comparables $\varepsilon = \{a,b\}$, $a < b$, et les composantes de $C(\varepsilon)$ sont $\mathcal L_{<a}$, $]a,b[$ et $\mathcal L_{>b}$ (\lot{156-3}{1}{8}). Les pages~11 à~13 montrent que l'intermédiarité stricte $\mathcal R$ et la comparabilité déterminent l'ordre à l'opposé près (\lot{156-3}{1}{11}, \lot{156-3}{1}{13}). La page~14 pose le problème : \q{la connaissance de $\mathrm{Drap}_2(\mathcal{L})$, et [\dots] de la décomp.\ de $C(\varepsilon)$ [\dots] en composantes connexes, implique la connaissance de la relation $\mathcal{R}(a,b,c)$} (\lot{156-3}{1}{14}) ; la page~15 calcule les trois composantes attachées à $a<b<c$, $X_{\{a,b\}} = \mathcal L_{>b}$, $X_{\{b,c\}} = \mathcal L_{<b}$, $X_{\{c,a\}} = \,]a,c[$, trouve la première intersection vide et les deux autres non vides, \q{divisibilité !}, et conclut \q{On gagne !} (\lot{156-3}{1}{15}).

\noindent\emph{Conventions.} Pour $a, b, c$ on note $\mathcal R(a,b,c)$ si $a < b < c$ ou $c < b < a$, et $\overline{\mathcal R}(a,b,c)$ si $a \leq b \leq c$ ou $c \leq b \leq a$. Pour $a \neq b$ comparables, $C(\{a,b\})$ est l'ensemble des éléments distincts de $a$ et $b$ et comparables aux deux ; deux de ses éléments sont dans le même \emph{morceau} s'ils sont tous deux au-dessous de $a$ et $b$, tous deux strictement entre eux, ou tous deux au-dessus : la relation ne dépend que de la paire. On note $X_{\{a,b\}}(t)$ le morceau de $C(\{a,b\})$ qui contient $t$. La donnée de la page~14 est celle de la comparabilité et, pour chaque paire comparable, de la partition de $C(\varepsilon)$ en morceaux, sans étiquette ; une bijection $e$ entre deux ensembles ordonnés \emph{respecte les découpages} si elle respecte la comparabilité et si, pour toute paire comparable $\{a,b\}$ et tous $s, t \in C(\{a,b\})$, $s$ et $t$ sont dans le même morceau de $C(\{a,b\})$ si et seulement si $e(s)$ et $e(t)$ sont dans le même morceau de $C(\{e(a), e(b)\})$.

\begin{enonce}[Page 15]
Soient $a < b < c$ dans un ensemble ordonné. Alors $X_{\{a,b\}}(c) = \mathcal L_{>b}$, $X_{\{b,c\}}(a) = \mathcal L_{<b}$ et $X_{\{c,a\}}(b) = \,]a,c[$. Si l'ordre est dense et si $p, q, r$ sont deux à deux distincts et comparables, $\mathcal R(p,q,r)$ a lieu si et seulement si $X_{\{p,q\}}(r) \cap X_{\{q,r\}}(p) = \emptyset$.
\end{enonce}

\begin{proof}
Si $b < t$ et $a < b$, un élément $s$ est dans le même morceau que $t$ exactement quand il est au-dessus de $a$ et de $b$ : $t$ n'est ni sous $a$ ni strictement entre $a$ et $b$. Un tel $s$ est dans $C(\{a,b\})$, d'où $X_{\{a,b\}}(c) = \mathcal L_{>b}$ ; de même $X_{\{b,c\}}(a) = \mathcal L_{<b}$, et, la relation étant symétrique en la paire, $X_{\{c,a\}}(b) = \,]a,c[$. Trois éléments deux à deux distincts et comparables forment une chaîne, et l'un des six ordres possibles a lieu. Si $q$ est au milieu, les deux morceaux sont $\mathcal L_{>q}$ et $\mathcal L_{<q}$, disjoints. Si le milieu est $p$, disons $q < p < r$, les morceaux sont $\mathcal L_{>p}$ et $]q,r[$, qui contiennent tous deux un point de $]p,r[$, non vide par densité ; les quatre autres cas sont semblables.
\end{proof}

\begin{enonce}[Pages 11 à 13]
Soient deux ensembles ordonnés filtrants croissants et décroissants, et $e$ une bijection de l'un sur l'autre qui respecte la comparabilité et l'intermédiarité stricte. Alors $e$ est un isomorphisme d'ordres ou un anti-isomorphisme. Plus précisément, si $a < b$, on a $u \leq v$ si et seulement s'il existe $x$ et $y$, avec $x = a$ ou $\mathcal R(b,a,x)$, $y = b$ ou $\mathcal R(a,b,y)$, tels que $\overline{\mathcal R}(x,u,y)$, $\overline{\mathcal R}(x,v,y)$ et $\overline{\mathcal R}(x,u,v)$. Enfin, $\overline{\mathcal R}(a,b,c)$ équivaut à : $\mathcal R(a,b,c)$, ou $b \in \{a,c\}$ avec $a$ et $c$ comparables.
\end{enonce}

\begin{proof}
La condition sur $x$ dit $x \leq a$, celle sur $y$ dit $y \geq b$. Si $u \leq v$, la filtration b) donne $x \leq u, a$ et $y \geq v, b$, et $x \leq u \leq v \leq y$ donne les trois conditions. Inversement, $x \leq a < b \leq y$ donne $x < y$ ; $\overline{\mathcal R}(x,u,y)$ ne peut alors avoir lieu que sous la forme $x \leq u \leq y$, et de même pour $v$ ; si $\overline{\mathcal R}(x,u,v)$ a lieu sous la forme $v \leq u \leq x$, alors $u = x$ et $v \leq x \leq v$, donc $u = v$ ; dans tous les cas $u \leq v$. La forme exacte de $\overline{\mathcal R}$ se vérifie cas par cas. Ainsi $\leq$ s'écrit à partir de $a$, $b$, de $\mathcal R$ et de la comparabilité seuls. Si l'ordre de départ n'a aucun couple $a < b$, c'est l'égalité, et la comparabilité respectée force l'égalité à l'arrivée. Sinon, soit $a < b$ ; $e(a)$ et $e(b)$ sont comparables et distincts. Si $e(a) < e(b)$, la description précédente, transportée par $e$, donne $u \leq v \Leftrightarrow e(u) \leq e(v)$. Si $e(b) < e(a)$, on remplace l'ordre d'arrivée par l'ordre opposé, qui a la même intermédiarité, la même comparabilité, et reste filtrant dans les deux sens.
\end{proof}

\begin{enonce}[Pages 11 à 15, \emph{l'ordre se lit sur les découpages}\piste{156-3-betweenness-from-cuts}]
Une bijection entre deux tronçons ordonnés qui respecte la comparabilité et, pour chaque paire comparable $\varepsilon$, la partition de $C(\varepsilon)$ en morceaux, est un isomorphisme d'ordres ou un anti-isomorphisme. La comparabilité seule n'y suffit pas : sur le tronçon $\mathbb Q$, la transposition de $0$ et $1$ respecte la comparabilité sans être ni croissante ni décroissante.
\end{enonce}

\begin{proof}
D'après l'énoncé de la page~15, appliqué dans les deux tronçons, et parce que $e$ envoie le morceau $X_{\{p,q\}}(r)$ sur le morceau $X_{\{e(p),e(q)\}}(e(r))$, la bijection $e$ respecte l'intermédiarité stricte ; on conclut par l'énoncé des pages~11 à~13. Pour $\mathbb Q$, les conditions a) à d) sont immédiates (minimum, maximum, milieu) ; tout couple y est comparable, et la transposition $\sigma$ de $0$ et $1$ vérifie $\sigma(1) = 0 < 1 = \sigma(0)$, alors que $\sigma(1) = 0 < 2 = \sigma(2)$.
\end{proof}

\begin{remarque}
Des conditions de la page~8, seules la divisibilité c) et la filtration b) servent : a), a$'$) et d) n'interviennent pas. La divisibilité sert à rendre non vides les deux autres intersections de la page~15, la filtration b) seulement dans le lemme des pages~11 à~13, et l'on n'a pas besoin de déduire $\mathrm{Drap}_2$ de $\mathcal R$, puisque la comparabilité fait partie des données de la page~14. L'énoncé vaut ainsi pour tout ensemble ordonné dense et filtrant dans les deux sens.
\end{remarque}

\begin{remarque}
La version large de l'intermédiarité est $\mathcal R(a,b,c)$, ou $b \in \{a,c\}$ \emph{avec $a$ et $c$ comparables}. La forme \q{$\mathcal R(a,b,c)$ ou $b \in \{a,c\}$} que la lecture substitue à la marge de la page~11 est fausse pour $a$ et $c$ incomparables : dans $\mathbb N^2$ ordonné produit, $\overline{\mathcal R}((0,1),(0,1),(1,0))$ n'a pas lieu. Le lemme de la page~13 ne rencontre pas ce cas, puisqu'il est appliqué avec $x < y$.
\end{remarque}

\noindent Ne sont pas démontrés ici : les axiomes des réseaux par découpages, que les pages~1 à~5 annoncent sans les écrire ; les remarques de la page~16 (la plus grande classe stable, et la question des conditions que doivent satisfaire $\mathcal R$ ou les $R_\varepsilon$) ; les parties modérées fermées et localement fermées (pages~17 à~21) ; la topologie d'un tronçon (pages~21 à~24) ; le bord, les tronçons et pseudo-tronçons et leur ordre (pages~24 à~28) ; l'algèbre des parties modérées et sa relativisation (pages~29 à~39).
